% =========================================================
% Tao--Mendelson Comparison
% =========================================================

\begin{remark*}[Exposition]
Tao and Mendelson construct the same integer system, but they choose different
entry points. Tao begins with natural-number pairs written as placeholder
differences \(a {-\!\!-} b\). The notation deliberately resembles
subtraction, even though subtraction has not yet been defined. Mendelson begins
with positive-integer pairs \((n,j)\) and explicitly takes equivalence classes.
The two constructions have the same quotient shape: a pair represents a formal
difference, and cross-addition determines when two representatives describe the
same integer.
\end{remark*}

\begin{remark*}[Exposition]
The pedagogical effect is different. Tao's notation keeps the computational
meaning visible. Addition, multiplication, negation, and order are introduced
by formulas that look like the formulas one expects for genuine differences.
The main burden is well-definedness: each formula must respect the equivalence
relation on representatives. Mendelson's notation makes the set-theoretic
quotient more explicit. The integer is an equivalence class from the start, so
the construction makes clear where representative choices enter.
\end{remark*}

\begin{remark*}[Exposition]
Their algebra also differs in style. Tao proves the algebra of
\(\mathbb{Z}\) concretely. After the operations are defined, the ring laws,
subtraction, no zero divisors, cancellation, and order are verified for the
constructed integer system itself. Mendelson pauses to abstract the pattern:
he defines rings, integral domains, ordered integral domains, positivity sets,
and cancellation principles before applying those ideas back to
\(\mathbb{Z}\). Tao's route is shorter and computational; Mendelson's route is
more reusable and structural.
\end{remark*}

\begin{remark*}[Exposition]
This explains why the two accounts are useful side by side. Tao is the better
guide for seeing why the integer operations have their familiar formulas.
Mendelson is the better guide for seeing which parts of the construction are
instances of general algebraic facts. The concrete account answers, ``How do
these integer formulas work?'' The abstract account answers, ``Which algebraic
structure has been built, and which later arguments can be reused in any
similar structure?''
\end{remark*}
